\documentclass[10pt,a4paper]{amsart}
\usepackage[margin=19mm]{geometry}
\usepackage{amsmath,amssymb,mathtools}
\usepackage[scr=rsfs,cal=euler]{mathalfa}
\usepackage{xfrac}
\usepackage[unicode,hidelinks,bookmarks=false]{hyperref}
\newcommand{\ie}{i.e.\ }
\newcommand{\cf}{cf.\ }
\newcommand{\resp}{respectively\ }
\newcommand{\Subsection}{\subsection}
\providecommand{\nobreakdash}{\nobreak\hbox{-}}
\setlength{\emergencystretch}{2em}
\multlinegap0pt
\usepackage{amssymb}
\usepackage{mathtools}
\newtheorem{proposition}{Proposition}
\newtheorem{lemma}{Lemma}
\newtheorem{theorem}{Theorem}
\newcommand{\AC}{\mathsf{AC}}
\newcommand{\HS}{\mathsf{HS}}
\newcommand{\Ord}{\mathrm{Ord}}
\newcommand{\VP}{\mathsf{VP}}
\newcommand{\ZF}{\mathsf{ZF}}
\newcommand{\forcesS}{\Vdash_{\mathrm{sym}}}
\newcommand{\namerank}{\operatorname{rank}_{\mathrm{name}}}
\newcommand{\rank}{\operatorname{rank}}
\newcommand{\restr}{\mathbin{\upharpoonright}}
\newcommand{\tc}{\operatorname{tc}}
\title{Vop\v{e}nka's Principle without Choice: Preservation under Symmetric Extensions}
\author{Tom de Groot, Wojciech Aleksander Wo{\l}oszyn}
\date{Source: arXiv:2609.06856v1 (CC BY 4.0)}
\begin{document}
\maketitle

\noindent\emph{Source and licence.} Vop\v{e}nka's Principle without Choice: Preservation under Symmetric Extensions. Version arXiv:2609.06856v1 (CC BY 4.0), distributed by arXiv under the Creative Commons Attribution 4.0 International licence, http://creativecommons.org/licenses/by/4.0/, as recorded in the arXiv licence metadata for this exact version and shown on the abstract page https://arxiv.org/abs/2609.06856v1. Full licence notice: \texttt{licenses/arxiv-2609.06856-CC-BY-4.0.txt}.\par\smallskip

\noindent\textit{Editorial scope.} Counted unit: the article's theorem thm:preservation (source0.tex lines 268--274). The article's complete original proof of it is delivered with it (source0.tex lines 276--340, one \texttt{proof} environment of 65 lines). No mathematics is written, repaired or re-derived: the statement and the proof are the article's own lines, in the article's own notation, and no retained line is substituted or re-flowed.  Retained uncounted as the source-local context the proof needs: lines 129--132; lines 182--188; lines 246--249.  Nothing was omitted from any retained span, so the package carries no omission record.  No retained line contains a citation, so the wrapper carries no bibliography and no bibliography entry.  No retained line contains a figure environment, an image-inclusion command or drawing code, so no image is delivered and the package declares no asset dependency.  Editorial formatting only, in addition: an AMS \texttt{amsart} preamble that restates the article's own declarations and macros used by the retained lines, namely the environments \texttt{proposition}, \texttt{lemma}, \texttt{theorem} on separate counters, together with the standard packages those lines need.  The retained environments print in this document as Proposition 1, Lemma 1, Theorem 1 in order of appearance, whereas the article prints the same items with the numbering of its own full document.  The complete native source of the article as distributed in the arXiv e-print for this exact version is bundled unchanged as \texttt{sources/209579/source0.tex}.
\begin{proposition}\label{prop:symmetric-ZF}
Suppose that $V\models\ZF$, that $\mathscr S$ is set-sized, and that $G\subseteq\mathbb P$ is $V$-generic.  Then
$W=\HS^G=\{\sigma^G:\sigma\in\HS\}$ is a transitive model of $\ZF$ with $V\subseteq W\subseteq V[G]$.
\end{proposition}

\begin{proposition}\label{prop:symmetric-lift}
Suppose that $V_\theta,V_\eta\models\ZF$, and let $j:V_\theta\to V_\eta$ be elementary.  Suppose that $j(z)=z$ for every $z\in\tc(\{\mathscr S\})$.  If $G\subseteq\mathbb P$ is $V$-generic and $W=\HS^G$, then there is a function $E_j\in W$ such that
\[
E_j(\sigma^G)=j(\sigma)^G\qquad(\sigma\in\HS\cap V_\theta).
\]
Thus, in $V[G]$, $E_j=j_G\restr\{\sigma^G:\sigma\in\HS\cap V_\theta\}$, where $j_G:V_\theta[G]\to V_\eta[G]$ is the ordinary forcing lift of $j$, defined by $j_G(\sigma^G)=j(\sigma)^G$.
\end{proposition}

\begin{lemma}\label{lem:ZF-ranks}

Assume $\ZF+\VP$.  Above every ordinal $\gamma$ there is a limit ordinal $\kappa$ such that no map from any $a\in V_\kappa$ is cofinal in $\kappa$.  Consequently $V_\kappa\models\ZF$.  Under $\AC$, such a $\kappa$ may be chosen measurable.
\end{lemma}

\begin{theorem}\label{thm:preservation}

Let $V\models\ZF+\VP$, with $\VP$ understood in the arbitrary-language sense fixed in the Introduction.  Let $\mathscr S=\langle\mathbb P,\Gamma,\mathcal F\rangle\in V$ be a set-sized symmetric system, and let $G\subseteq\mathbb P$ be $V$-generic.  Then
\[
W=\HS^G\models\ZF+\VP.
\]
\end{theorem}

\begin{proof}

Proposition~\ref{prop:symmetric-ZF} gives $W\models\ZF$.  Fix a formula $\varphi(x,y)$, a parameter $a\in W$, and a set-sized language $L\in W$ such that
$\mathcal A=\{M\in W:W\models\varphi(M,a)\}$ is a proper class of $L$-structures.  Choose $\dot a,\dot L\in\HS$ with $\dot a^G=a$ and $\dot L^G=L$.  For $p\in\mathbb P$, let $R_p$ be the class of ordinals $\alpha$ for which there is a name $\tau\in\HS$ satisfying
\begin{equation}\label{eq:rank-witness}
p\forcesS ``\tau\text{ is an }\dot L\text{-structure},\ \varphi(\tau,\dot a),
\text{ and }\rank(\tau)=\check\alpha.''
\end{equation}

{In \eqref{eq:rank-witness}, $\rank(\tau)$ denotes the rank of the value of $\tau$ in the generic extension.  By contrast, $\namerank(\tau)$ is the von Neumann rank of the name $\tau$ itself.}

There is a condition $q\in G$ for which $R_q$ is unbounded in $\Ord$.  Suppose not.  Let $D_{\mathrm{bd}}=\{p\in\mathbb P:R_p\text{ is bounded}\}$ and, for $p\in D_{\mathrm{bd}}$, let $b(p)$ be the least ordinal such that $R_p\subseteq b(p)$.  The domain $D_{\mathrm{bd}}$ is a subset of the set $\mathbb P$, and definability of symmetric forcing together with Replacement makes $b$ a set function in $V$.  Our supposition gives $G\subseteq D_{\mathrm{bd}}$.

In $V[G]$, set $\delta=\sup b``G$.  This is an ordinal, so preservation of ordinals gives $\delta\in V\subseteq W$.  If $M\in\mathcal A$, choose $\tau\in\HS$ with $\tau^G=M$ and let $\alpha=\rank(M)$.  By the symmetric truth lemma, some $p\in G$ forces the assertion in \eqref{eq:rank-witness}; hence $\alpha<b(p)\leq\delta$.  Thus $\mathcal A\subseteq V_\delta^W$, and Separation in $W$ makes $\mathcal A$ a set, a contradiction.

For $\alpha\in R_q$, let $\Psi_\alpha(\tau)$ assert that $\tau\in\HS$ and that \eqref{eq:rank-witness} holds with $q$ in place of $p$.  We use the following definitions:
\[
\nu_\alpha=\min\{\namerank(\tau):\Psi_\alpha(\tau)\},
\qquad
X_\alpha=\{\tau\in V_{\nu_\alpha+1}:\namerank(\tau)=\nu_\alpha
\text{ and }\Psi_\alpha(\tau)\}.
\]
{The ordinal $\nu_\alpha$ exists, $X_\alpha$ is a nonempty set by Separation, and both are uniformly definable from $\alpha$ and the fixed parameters.}

Choose $\rho>\omega$ such that $\tc(\{\mathscr S,\dot a,\dot L\})\subseteq V_\rho$.  By Lemma~\ref{lem:ZF-ranks}, for each $\alpha\in R_q$ there are arbitrarily high $\theta$ with $V_\theta\models\ZF$.  The relevant rank is
\[
\theta_\alpha=\min\bigl\{\theta:V_\theta\models\ZF\text{ and }
\theta>\max\{\rho,\alpha,\rank(X_\alpha)\}+\omega\bigr\}.
\]
In the common language consisting of a binary relation symbol, two distinguished constant symbols, and constant symbols $c_z$ for $z\in V_\rho$, use the structure
\[
\mathcal M_\alpha=
\bigl\langle V_{\theta_\alpha},\in,\alpha,X_\alpha,(z)_{z\in V_\rho}\bigr\rangle
\qquad(\alpha\in R_q).
\]

The structures $\mathcal M_\alpha$ form a definable proper class since $R_q$ is unbounded.  By $\VP$, there is an elementary embedding $j:\mathcal M_\alpha\to\mathcal M_\beta$ with $\alpha\ne\beta$.  The distinguished constants give $j(\alpha)=\beta$, $j(X_\alpha)=X_\beta$, and $j(z)=z$ for every $z\in V_\rho$.

Fix $\tau\in X_\alpha$.  Then $j(\tau)\in X_\beta$, and the definitions of $X_\alpha$ and $X_\beta$ give
\begin{equation}\label{eq:q-witness}
q\forcesS ``\tau\text{ and }j(\tau)\text{ are }\dot L\text{-structures satisfying }
\varphi(-,\dot a)\text{ with ranks }\check\alpha\text{ and }\check\beta.''
\end{equation}
Since $j$ fixes $V_\rho$ pointwise, it fixes every member of $\tc(\{\mathscr S,\dot a,\dot L\})$.

Apply Proposition~\ref{prop:symmetric-lift} to obtain the elementary lift
$j_G:V_{\theta_\alpha}[G]\to V_{\theta_\beta}[G]$ and a function $E_j\in W$ agreeing with $j_G$ on values of names in $\HS\cap V_{\theta_\alpha}$.  Put $M=\tau^G$ and $N=j(\tau)^G$.  Then $j_G(M)=N$.

If $s\in L=\dot L^G$, then $s=\sigma^G$ for some $\langle\sigma,p\rangle\in\dot L$ with $p\in G$.  Since $\sigma\in\operatorname{cl}(\dot L)\subseteq V_\rho$, we have $j(\sigma)=\sigma$, and hence $j_G(s)=s$.  Thus $j_G$ fixes $L$ pointwise.  Moreover, $j(\dot L)=\dot L$ and $j_G(L)=L$.

By the coding convention of the Introduction, $\lvert M\rvert\subseteq\tc(M)$.  An induction on name evaluation shows that every element of $\tc(M)$ is $\sigma^G$ for some $\sigma\in\operatorname{cl}(\tau)$.  Since $\tau\in\HS\cap V_{\theta_\alpha}$, every such $\sigma$ belongs to $\HS\cap V_{\theta_\alpha}$.  Hence $|M|\subseteq\operatorname{dom}(E_j)$, and $e=E_j\restr|M|$ belongs to $W$.  In $V[G]$, the map $e$ is $j_G\restr|M|$.  Since the universe of a structure is uniformly definable from its fixed set code,
\[
j_G(|M|)=|j_G(M)|=|N|,
\qquad j_G``|M|\subseteq|N|.
\]

{Satisfaction for set-sized structures is absolute between the transitive models under consideration.  Since $j_G$ fixes $L$ pointwise, elementarity gives}
\[
M\models\psi(\bar m)\quad\text{if and only if}\quad
N\models\psi(e(\bar m))
\]
{for every $L$-formula $\psi$ and every finite tuple $\bar m$ from $|M|$.  Thus $W$ regards $e$ as an $L$-elementary embedding from $M$ to $N$.}

Finally, $q\in G$ and \eqref{eq:q-witness} imply that $M,N\in\mathcal A$, with ranks $\alpha$ and $\beta$.  Thus $M\neq N$.  This proves the arbitrary instance of $\VP$, so $W\models\ZF+\VP$.
\end{proof}

\end{document}
